\section{Related Work and Research Gap}\label{sec:related-work}

\subsection{Earth observation for archaeological and cultural-heritage landscapes}

Earth observation in archaeology has developed from site prospection towards landscape interpretation, change documentation and management support. Optical, radar, thermal and elevation data can reveal crop or soil marks, topographic traces, settlement patterns and changing environmental context across areas that would be costly to survey repeatedly \citep{lasaponara2011satellite,tapete2018geosciences,opitz2018trends,yao2023overview}. The shift is important for conservation because a distributed archaeological property is affected not only by processes at visible masonry but also by changes in its surroundings. Studies at Paphos, for example, used satellite imagery and GIS to examine urban expansion and clusters of natural and anthropogenic pressure \citep{agapiou2015paphos,agapiou2016paphos}. At Angkor, remote sensing and hydrological modelling were used to map flood hazard within a World Heritage setting \citep{liu2019angkor}. The unit of analysis consequently ranges from an archaeological feature to an entire landscape, and the meaning of accuracy changes with it. Detecting a feature, mapping land cover and prioritising an inspection are related but distinct tasks. These applications show that synoptic observation can organise field attention, while their different targets also show why ``remote-sensing assessment'' is not one uniform evidential claim.

Recent work has expanded the scale and automation of archaeological screening. Machine learning has supported object detection from airborne laser scanning and multisensor imagery, while space-based platforms have been proposed as part of wider heritage safeguarding systems \citep{lambers2019ml,orengo2020mounds,argyrou2022review,luo2023space}. Automation improves consistency and coverage when the target class is adequately represented, but it does not remove dependence on training labels, spatial resolution and field verification. A classifier trained to detect mounds answers a different question from one trained on land-cover labels, and neither becomes a condition-assessment model by being applied near a monument. Reviews of archaeological remote sensing repeatedly identify an imbalance between technical capability and interpretive validation \citep{opitz2018trends,argyrou2022review}. For CHIP, the methodological consequence is to use automated computation for repeatable screening while preserving a human-inspection endpoint. The system ranks questions for field teams; it does not automate a diagnosis of heritage condition.

\subsection{Multi-temporal indicators, land-surface interpretation and the condition boundary}

The Landsat archive makes long-term comparison possible through an unusually consistent record of moderate-resolution land observation \citep{hansen2012landsat,wulder2019landsat}. Its spatial resolution is suitable for neighbourhood-scale landscape monitoring, but individual pixels can contain mixtures of archaeological fabric, soil, vegetation, paths and buildings. Sensor transitions and atmospheric or seasonal differences further complicate temporal comparison. Cross-sensor work has documented systematic differences between Landsat 7 and Landsat 8 reflectance and vegetation indices, supporting harmonisation and cautious interpretation rather than unqualified pooling \citep{roy2016landsat}. Multi-date composites can improve usable coverage and suppress transient contamination, yet they represent intervals rather than exact dates of change. A median composite can reduce the influence of outliers while also hiding short disturbances that occur between or within the selected windows. It therefore trades event timing for spatial completeness and comparability. CHIP consequently uses fixed seasonal windows, quality masking, common analytical support and epoch terminology rather than implying continuous annual observation.

Spectral indices provide compact descriptions of surface response, but their names can encourage interpretations stronger than their physics. NDVI contrasts near-infrared and red reflectance and is widely used to characterise green vegetation \citep{tucker1979ndvi}. MNDWI contrasts green and short-wave infrared reflectance to enhance open-water expression \citep{xu2006mndwi}, while NDBI was designed to highlight built-up surfaces using short-wave infrared and near-infrared response \citep{zha2003ndbi}. Bare soil, sparse vegetation, moisture state, material mixtures and sensor characteristics can produce overlapping signals. At archaeological sites, a lower vegetation index may reflect removal, senescence, soil background or altered management; a higher NDBI may reflect construction or exposed ground; and a wetness-index change may reflect water, damp soil or mixed land cover. Directional transformation into a ``pressure'' variable adds decision meaning but does not resolve that equifinality. Threshold convergence is more informative than one index alone when several surface dimensions agree, yet the thresholds still require empirical and field scrutiny. The indices are therefore evidence about relative landscape context, not direct observations of stone, mortar or structural integrity.

Scale adds a second interpretive constraint. The statistical properties of remotely sensed scenes change with grain and support, and an object that dominates a small neighbourhood can be diluted or combined with different surroundings at a larger radius \citep{woodcock1987scale}. Heritage studies often select a buffer distance for convenience or comparability, but the resulting summary can be sensitive to that decision. The selected radius determines the mixture of archaeological space, roads, vegetation, settlement and terrain entering each component statistic. Point-centred circles are especially vulnerable to semantic overreach because their geometric precision can resemble an official boundary even when none was supplied. Scale comparison is therefore both a numerical sensitivity test and an interpretation test. CHIP treats 250, 500 and 1,000~m radii as alternative analytical supports, compares their rankings explicitly and never presents them as legal buffers or property extents.

\subsection{Climate forcing, hydro-geomorphology and archaeological-landscape susceptibility}

Climate-related heritage research documents pathways involving temperature, rainfall, humidity, drought, flooding, erosion, biological activity and interacting material responses \citep{fatoric2017threatened,sesana2021climate,orr2021climate}. Evidence strength varies substantially. Some studies use measured fabric responses or documented loss, while others map exposure or project climatic suitability for damaging processes. The word ``climate'' can therefore refer to a long-term change in forcing, a modelled future scenario, an observed extreme event or the environmental context of one survey. Reviews caution against moving directly from any of these levels to a local deterioration claim. Landscape-based monitoring within the World Heritage system has been proposed as a way to connect climate observation with management, but it still requires information about the heritage attributes and mechanisms at issue \citep{guzman2020monitoring}. CHIP adopts that narrower role: climate data describe forcing and acquisition context, not observed monument response.

Case studies demonstrate why both long-term climate and event-specific context matter. Reimann et al. assessed future coastal flooding and erosion exposure for Mediterranean World Heritage sites, while Hollesen et al. synthesised rapidly changing preservation conditions across Arctic archaeological archives \citep{reimann2018mediterranean,hollesen2018arctic}. These studies have explicit hazard or preservation mechanisms and should not be treated as interchangeable with an inland, point-centred landscape screen. A trend in annual precipitation can coexist with an unusually wet or dry acquisition window, and two gridded products can disagree because they assimilate observations and represent precipitation differently. At Taxila, the available weather record is gridded and lacks local station validation. Its defensible contribution is to describe annual and seasonal variability, trends, product disagreement and antecedent conditions around satellite acquisition. The number of independent image epochs remains too small for strong attribution, so climate-spectral relationships are exploratory associations.

Terrain supplies a more spatially differentiated but still indirect form of evidence. Elevation gradients can be transformed into slope, flow direction, contributing area, topographic position and wetness proxies. These variables structure how surface and near-surface water may converge across a landscape \citep{beven1979topmodel,moore1991terrain,quinn1991flow,tarboton1997flow}. Their physical rationale is stronger than an arbitrary proximity score, but their meaning depends on elevation resolution, routing assumptions and unobserved infrastructure. A drainage-proximity layer derived from a digital elevation model does not show a functioning channel, blocked culvert, flood depth or erosion rate. Heritage flood and multi-hazard studies that possess event, exposure or vulnerability information can estimate more specific outcomes \citep{liu2019angkor,guerriero2022derwent,arrighi2023flood}. CHIP instead labels terrain and drainage variables as susceptibility descriptors and uses them to define what should be checked on site.

\subsection{Transparent composite prioritisation, scale, correlation and uncertainty}

GIS-based multi-criteria decision analysis offers a practical means of combining heterogeneous evidence through normalisation, weighting and aggregation \citep{malczewski2006gis}. Heritage applications have used such architectures to integrate natural and anthropogenic pressures, flood context and multi-hazard susceptibility \citep{hadjimitsis2013cyprus,guerriero2022derwent,yang2023yongtai}. Their strength is legibility: a map can connect multiple observations to a decision objective. Their weakness is that normalisation choices, correlated indicators, weights and aggregation rules can become hidden behind one polished score. Compensatory aggregation can also allow a low value in one domain to be offset by a high value in another, which may or may not match the conservation decision. A value reported to several decimals may therefore remain conditional on judgement and data support. Transparency requires the score architecture and its alternatives to be published together. Sensitivity analysis converts those choices into results that can be examined. Spatial weight analysis has shown that alternative criterion weights can alter both local suitability values and their ordering \citep{chen2010sensitivity,feizizadeh2014uncertainty}. Factor ablation tests whether the conclusion depends on one variable, while domain ablation asks whether one evidence family controls the outcome. Correlation analysis identifies indicators that may duplicate the same ordering information. These procedures are particularly relevant when several spectral transformations respond to a shared surface-cover gradient. CHIP groups correlated indicators hierarchically, removes factors and domains in turn, and perturbs domain weights rather than treating the declared weighting as uniquely correct.

Uncertainty in a priority map also has a spatial component. Classical resampling provides a general route to empirical uncertainty estimation \citep{efron1979bootstrap}, but independent pixel resampling is poorly matched to spatially autocorrelated landscapes. Block-based resampling retains part of the local dependence and asks whether ranks persist when different spatial units contribute to the component summaries. This does not capture every source of uncertainty. It does, however, prevent the fixed ranking from being interpreted as the only plausible ordering. The methodological consequence for CHIP is to report rank distributions and interval overlap alongside the fixed 500~m score, so that components with indistinguishable ranking support can be placed in the same inspection tier.

\subsection{Spatial validation, proxy labels and reproducible geospatial science}

Spatial prediction requires validation that matches the intended transfer. Randomly splitting neighbouring pixels can place strongly related observations in both training and test data, creating an optimistic estimate of performance away from the sampled locations. Spatially structured cross-validation and buffered or blocked holdouts reduce that leakage, although the correct design depends on whether the intended claim concerns interpolation, extrapolation or map accuracy \citep{roberts2017cv,valavi2019blockcv,ploton2020spatial,wadoux2021spatial}. The distinction is essential here because the model is evaluated on a geographically separated outer region. Its score addresses transfer to that region under one partition, not universal transfer across Taxila or beyond it. The response label sets a separate boundary. Accuracy assessment requires a reference design appropriate to the thematic claim, including attention to sampling, response quality and analysis \citep{olofsson2014accuracy,stehman2019accuracy}. ESA WorldCover offers globally consistent 10~m land-cover classes derived from Sentinel-1 and Sentinel-2, but it remains a mapped product with class-specific error rather than field truth \citep{zanaga2021worldcover,esa_worldcover}. Agreement can reflect genuine spectral separability, shared spatial structure, inherited label error or a combination of all three. A model trained to reproduce WorldCover can therefore reveal whether Landsat features carry geographically transferable information aligned with that land-cover proxy. It cannot independently validate WorldCover, archaeological condition or the CHIP ranking. Class-wise confusion, calibration and block-bootstrap intervals qualify model performance, but no metric can enlarge the meaning of the response labels. CHIP therefore describes every model metric as WorldCover-derived proxy land-cover agreement.

Computational transparency is necessary because geospatial results depend on source versions, projections, masks, resampling rules, random seeds and software environments. Reproducibility studies in GIScience have shown that narrative method descriptions alone often do not permit an analysis to be regenerated \citep{nust2018reproducible}. General scientific-computing guidance similarly recommends preserving data organisation, scripts, environments and intermediate decisions \citep{peng2011reproducible,stodden2016enhancing,wilson2017practices}. For a composite priority, reproducibility must extend beyond the final raster to the transformations, weights, exclusions and uncertainty draws that produced it. CHIP follows this principle through source manifests, fixed configurations, code-linked outputs, validation records, checksums and a claim-evidence register. Reproducibility does not validate the conservation interpretation, but it makes the analytical chain open to correction, substitution and rerunning when field evidence becomes available.

\subsection{Taxila-specific evidence and the unresolved gap}

The official World Heritage record establishes Taxila as a serial property of multiple archaeological components, while the study by Khan et al. provides the closest empirical assessment of its natural and anthropogenic threats \citep{unesco_taxila,khan2022taxila}. Khan et al. combined remote sensing with ground-based seismic measurements and framed their method as a cost-conscious contribution to hazard documentation. The present study does not displace that evidence. It lacks equivalent field-condition and geophysical observations, and its analytical circles do not substitute for official component polygons. Its point of departure is the temporal and decision-uncertainty problem: how to compare five satellite epochs on common support, connect them to acquisition-window weather, integrate terrain and drainage context, and determine whether a component ranking survives alternative radii, factors, domains, weights and spatial samples. This distinction avoids constructing a research gap by understating the earlier Taxila study. The gap concerns the transparency and stability of an inspection-priority decision, not the existence of prior hazard evidence.

Evidence has established that Earth observation, GIS, climate records and terrain analysis can support heritage monitoring, and that Taxila faces a combination of natural and anthropogenic pressures. It has not adequately resolved how those evidence streams can be integrated at Taxila without conflating landscape pressure with monument damage, treating coarse weather as fine-scale hazard, double-counting correlated indicators, assuming one neighbourhood scale, or presenting a fixed priority rank without uncertainty. The unresolved issue is thus not a shortage of maps, but the traceability and stability of the decision made from them. This study addresses that problem through the multi-epoch CHIP design, explicit common support, climate-window matching, hierarchical indicator grouping, multiscale comparison, ablation, weight perturbation, spatial block-bootstrap ranking and geographically buffered proxy-label evaluation. It retains the central inference boundary that the output is a relative guide to field-inspection sequencing, not an independently validated condition, damage or legal-risk assessment. This framing keeps the decision operational while leaving condition diagnosis to appropriately designed field investigation.
